% !TeX program = xelatex
% UTF-8; compile twice with XeLaTeX. No external figures or input files are needed.
\IfFileExists{ctexart.cls}{
  \documentclass[UTF8,11pt,a4paper]{ctexart}
}{
  \documentclass[11pt,a4paper]{article}
  \usepackage{fontspec}
  \setmainfont[AutoFakeBold=2,AutoFakeSlant=0.15]{Noto Sans CJK SC}
  \XeTeXlinebreaklocale "zh"
  \XeTeXlinebreakskip=0pt plus 1pt
}
\usepackage[margin=24mm,headheight=16pt]{geometry}
\usepackage{amsmath,amssymb,amsthm,mathtools}
\usepackage{booktabs,array,longtable,fancyhdr}
\usepackage[hidelinks,bookmarksnumbered]{hyperref}
\hypersetup{pdftitle={XP-算子研究草稿：紧性、线性差、自然性与Cauchy相关集},pdfauthor={XP-算子项目}}
\theoremstyle{definition}
\newtheorem{theorem}{定理}[section]
\newtheorem{proposition}[theorem]{命题}
\newtheorem{lemma}[theorem]{引理}
\newtheorem{corollary}[theorem]{推论}
\newtheorem{definition}[theorem]{定义}
\newtheorem{example}[theorem]{例}
\newtheorem{remark}[theorem]{注}
\renewcommand{\proofname}{\normalfont\bfseries 证明}
\renewcommand{\contentsname}{目录}
\newcommand{\K}{\mathbb K}
\newcommand{\Q}{\mathbb Q}
\newcommand{\R}{\mathbb R}
\newcommand{\C}{\mathbb C}
\newcommand{\Z}{\mathbb Z}
\newcommand{\N}{\mathbb N}
\newcommand{\Hom}{\operatorname{Hom}}
\newcommand{\End}{\operatorname{End}}
\newcommand{\im}{\operatorname{im}}
\newcommand{\coker}{\operatorname{coker}}
\newcommand{\Span}{\operatorname{span}}
\newcommand{\supp}{\operatorname{supp}}
\newcommand{\id}{\operatorname{id}}
\newcommand{\ev}{\operatorname{ev}}
\newcommand{\D}{\mathcal D}
\newcommand{\T}{\mathcal T}
\newcommand{\calA}{\mathcal A}
\newcommand{\calM}{\mathcal M}
\newcommand{\calN}{\mathcal N}
\newcommand{\Sol}{\operatorname{Sol}}
\DeclareMathOperator{\rank}{rank}
\setlength{\parskip}{0.3em}
\setlength{\emergencystretch}{3em}
\renewcommand{\arraystretch}{1.17}
\pagestyle{fancy}
\fancyhf{}
\fancyhead[L]{\small XP-算子研究草稿}
\fancyhead[R]{\small 2026年10月9日}
\fancyfoot[C]{\thepage}
\begin{document}
\begin{titlepage}
\vspace*{20mm}
\begin{center}
{\LARGE\bfseries XP-算子研究草稿\par}
\vspace{7mm}
{\Large 紧性、线性差与非齐次方程\par}
\vspace{3mm}
{\large 自然性及 Cauchy 相关集的补充研究\par}
\vspace{10mm}
{2026年10月9日修订\par}
\end{center}
\vspace{10mm}
\noindent\textbf{内容摘要。}
本稿整理本项目关于紧算子、线性差与自然性的研究。
本次修订采用2026年10月9日的术语：线性差专指对函数变量满足 $K$-线性的算子，仅可加的对象称为加性差。
首先给出一般紧算子的有限依赖、无交互项和离散乘积拓扑刻画；
其次把 $Q=P-h$ 的全域函数问题化为线性映射的纤维，并在标量值与有限维向量值情形给出完整的有限关系判据、有限相容性定理及解空间参数化。
无限维值域单独处理，给出模同态延拓判据和有限相容性失效的反例。
随后明确若干不同的自然性条件，证明平移等变算子的群代数表述，以及对全部线性映射自然的有限取值公式。
本次进一步补充一般标量线性差与紧线性差的差别、无限维值域的双对偶障碍、固定空间上的线性群与自同态自然性分类，以及自然模板的矩和一元伸缩方程。
最后整理加法 Cauchy 缺陷的非齐次解、差分多项式、复代数零集和实光滑性，并收录乘法 Cauchy 关系的近期结果。

\medskip
\noindent\textbf{整理原则。}
所有标为定理、命题或推论的结论均在文内给出证明。
自然性章节先明确新增假设，再作推导，不把不同自然性要求混为一谈。
本稿不重述既有整函数乘法缺陷分类长篇报告，也不以该报告中的分类作为本稿证明的前提。
第\ref{sec:scope}节说明与既有草稿的对应及适用边界。

\medskip
\noindent\textbf{约定。}
除另有说明，函数是任意集合映射，不附加连续性或可测性。
使用通常的集合论选择公理，因而允许选择向量空间基、补空间和线性延拓。
“紧”始终采用本项目的有限取值可分解定义，不指泛函分析中的紧算子。
\vfill
\noindent\small 源文件为单文件 UTF-8 文档，使用 XeLaTeX 连续编译两次即可生成目录与交叉引用。
\end{titlepage}
\begingroup
\small
\setlength{\parskip}{0pt}
\tableofcontents
\endgroup
\clearpage

\section{基本对象、问题类型与线性性的层次}
设 $R,S$ 为交换环，$V,W$ 分别为 $R,S$-模，$n\ge1$，并令
\[
 I=V,\qquad J=V^n,\qquad W^I=\{f:I\to W\}.
\]
对映射 $P:W^I\times J\to W$，定义
\[
 X_P(f)=\{v\in J:P(f,v)=0\}.
\]
它称为 $f$ 的 $P$-域，也沿用讨论中的 $P$-Area 一词；这里的“域”不指代数意义的域。

\begin{definition}[三类问题与本稿的非齐次问题]
$P$-函数问题要求 $X_P(f)=J$；反 $P$-函数问题要求 $X_P(f)=\varnothing$；
$P$-域问题要求对给定 $f$ 求 $X_P(f)$。
给定 $h:J\to W$，令 $Q(f,v)=P(f,v)-h(v)$。本稿所说的 $Q$-问题均指
\begin{equation}\label{eq:inhom}
 P(f,v)=h(v)\qquad(v\in J),
\end{equation}
即全域 $Q$-函数问题。要求 $X_Q(f)$ 恰为某个指定真子集，是另一个实现问题。
\end{definition}

\begin{definition}[加性差与线性差：修订后的约定]
若对全部 $f,g,v$ 有
\[
 P(f+g,v)=P(f,v)+P(g,v),
\]
则称 $P$ 为\textbf{加性差}。
当 $R=S=K$ 为域时，若此外对每个 $a\in K$ 还满足
\[
 P(af,v)=aP(f,v),
\]
则称 $P$ 为\textbf{线性差}，必要时强调为 $K$-线性差。
本文统一采用这一修订，不再用“线性差”指仅满足可加性的对象。
\end{definition}
可加性使 $P(0,v)=0$。特征零时，可加映射自动为 $\Q$-线性；
特征 $p$ 时自动为 $\mathbb F_p$-线性。一般并不自动为 $K$-线性。
例如 $K=\C$ 时 $P(f,v)=\overline{f(0)}$ 可加，却不是复线性的。

\begin{proposition}[一般非齐次方程的纤维]
若 $P$ 为 $K$-线性差，令
\[
 \T:W^I\to W^J,\qquad (\T f)(v)=P(f,v).
\]
则方程\eqref{eq:inhom}可解当且仅当 $h\in\im\T$；有一个解 $f_0$ 时，全部解为
\begin{equation}\label{eq:torsor}
 \Sol(h)=f_0+\ker\T.
\end{equation}
对某个可解的 $h$ 解唯一，当且仅当 $\ker\T=0$；这也等价于每个可解的右端均有唯一解。
每个右端恰有一解，当且仅当 $\T$ 为同构。
\end{proposition}
\begin{proof}
$\T f=h=\T f_0$ 当且仅当 $\T(f-f_0)=0$，得到\eqref{eq:torsor}。
存在性、唯一性和全体右端的结论分别就是像、核及双射的定义。
\end{proof}
若只假设可加性，同一证明在加法群层次成立，在素域层次给出仿射空间。
即便 $V,W$ 有限维，$W^V$ 通常仍是巨大的函数空间；不能对 $V,W$ 的有限维数直接套用一个有限矩阵的秩公式。

\section{紧算子的完整等价刻画}\label{sec:compact}
本节只需要 $W$ 是阿贝尔群，$I\ne\varnothing$。对 $x\in I,w\in W$，记
\[
 (\delta_xw)(y)=\begin{cases}w,&y=x,\\0,&y\ne x.\end{cases}
\]
\begin{definition}[项目中的紧性]
$P$ 紧，是指对每个 $v\in J$，存在有限集 $F_v\subseteq I$ 及映射 $U_{x,v}:W\to W$，使
\begin{equation}\label{eq:compact}
 P(f,v)=\sum_{x\in F_v}U_{x,v}(f(x))\qquad(f\in W^I).
\end{equation}
有限集与局部映射可以依赖 $v$；不要求一个统一的有限集合适用于全部 $v$。
\end{definition}

\begin{theorem}[紧性、有限依赖与无交互项]\label{thm:compact}
固定 $v$，记 $T(f)=P(f,v)$、$c=T(0)$。以下条件等价：
\begin{enumerate}
\item $T$ 具有\eqref{eq:compact}的形式。
\item 存在有限集 $F\subseteq I$，$T$ 只依赖 $f|_F$，且对支集不交的 $f,g$ 有
\begin{equation}\label{eq:orthadd}
 T(f+g)=T(f)+T(g)-c.
\end{equation}
\item $W$ 取离散拓扑、$W^I$ 取乘积拓扑时，$T$ 在 $0$ 处连续，且满足\eqref{eq:orthadd}。
\item 存在唯一的有限集 $E$ 和唯一的一族非零映射 $A_x:W\to W$，满足 $A_x(0)=0$，且
\begin{equation}\label{eq:canonical}
 T(f)=c+\sum_{x\in E}A_x(f(x)).
\end{equation}
\end{enumerate}
这里允许 $E=\varnothing$。规范局部映射为
\begin{equation}\label{eq:Ax}
 A_x(w)=T(\delta_xw)-c,
\end{equation}
并将恒零的项删去。对这类 $T$，在 $0$ 连续也等价于处处连续。
\end{theorem}
\begin{proof}
(1)$\Rightarrow$(2)：有限依赖显然。若 $\supp f\cap\supp g=\varnothing$，则每个坐标至少有一个取值为零，故
\[
 U_x(f(x)+g(x))=U_x(f(x))+U_x(g(x))-U_x(0).
\]
求和并用 $c=\sum_xU_x(0)$ 得\eqref{eq:orthadd}。

(2)$\Rightarrow$(4)：记 $f_F$ 为 $f$ 在 $F$ 上的截断。有限依赖给出 $T(f)=T(f_F)$。
对有限和 $f_F=\sum_{x\in F}\delta_xf(x)$ 逐次用\eqref{eq:orthadd}，得到
\[
 T(f)=c+\sum_{x\in F}\bigl(T(\delta_xf(x))-c\bigr).
\]
按\eqref{eq:Ax}定义 $A_x$，删去恒零项即得规范式。任一满足 $A_x(0)=0$ 的分解，在 $\delta_xw$ 处代入后都必须满足\eqref{eq:Ax}，故唯一。

(4)$\Rightarrow$(1)：若 $E\ne\varnothing$，把常数 $c$ 加入其中一项；若 $E=\varnothing$，任选 $x_0\in I$，令 $U_{x_0}\equiv c$ 即可。

(2)$\Rightarrow$(3)：固定有限个坐标的取值后输出不变，故 $T$ 处处连续。
反之，(3)意味着存在有限 $F$，只要 $g|_F=0$ 就有 $T(g)=c$。
将任意 $f$ 分解成支集不交的 $f_F+f_{I\setminus F}$，得到
\[
 T(f)=T(f_F)+T(f_{I\setminus F})-c=T(f_F).
\]
这就是有限依赖，再结合假设中的\eqref{eq:orthadd}得到(2)。
\end{proof}

\begin{remark}[不能删去无交互项条件]
有限依赖和连续性本身均不足以推出上述紧性。
例如 $W=\R$，取不同的 $a,b\in I$，令 $T(f)=f(a)f(b)$。
它只依赖两个坐标且连续，但
$T(\delta_a1+\delta_b1)=1$，而两项单独输入的输出都为零，因此不满足\eqref{eq:orthadd}。
\end{remark}

\begin{corollary}[加性差与线性差的紧性]\label{cor:lincompact}
若 $T:W^I\to W$ 可加，则以下条件等价：项目紧性；有限坐标依赖；离散乘积拓扑下在零点连续。
此时规范式中 $c=0$，且每个 $A_x$ 可加。若 $T$ 还是 $K$-线性的，则每个 $A_x\in\End_K(W)$。
\end{corollary}
\begin{proof}
可加性自动给出\eqref{eq:orthadd}，且 $T(0)=0$。
由 $A_x(w)=T(\delta_xw)$ 及 $\delta_x$ 的线性性立即得到相应的线性结论。
\end{proof}

\begin{corollary}[局部线性性无须另作假设]\label{cor:autolocal}
按修订后的定义，任何紧线性差都具有唯一的归一化表示
\[
 P(f,v)=\sum_{x\in F_v}T_{v,x}f(x),\qquad
 T_{v,x}\in\End_K(W),\quad T_{v,x}\ne0,
\]
其中 $F_v$ 有限，且 $T_{v,x}(w)=P(\delta_xw,v)$。
因此研究紧线性差时，无须再附加“存在一组线性的局部映射”这一假设。
\end{corollary}
\begin{proof}
规范分解的常数项为 $P(0,v)=0$；对 $w$ 的加性和齐次性均从 $P$ 继承。
唯一性来自定理\ref{thm:compact}。原始非归一化分解中的局部项可能带互相抵消的常数，故这里的结论针对归一化表示。
\end{proof}

\section{一般标量线性差：紧性之外的完整判据}\label{sec:noncompact}
本节只假设 $W=K$ 且 $P$ 为线性差，\textbf{暂不假设紧性}。
记 $E=K^I,F=K^J$，仍以 $\T:E\to F$ 表示 $P$，并令
\[
 \T^*:F^*\to E^*,\qquad \T^*(\Lambda)=\Lambda\circ\T,
\]
其中星号表示\textbf{完整代数对偶}。

\begin{theorem}[一般线性差的对偶存在唯一性]\label{thm:fullalgebraic}
方程 $\T f=h$ 可解，当且仅当
\begin{equation}\label{eq:fullcompat}
 \Lambda\in F^*,\quad \Lambda\circ\T=0
 \quad\Longrightarrow\quad \Lambda(h)=0.
\end{equation}
非空解集仍为 $f_0+\ker\T$。此外，每个右端的解都至多一个当且仅当 $\T^*$ 满射；每个右端可解当且仅当 $\T^*$ 单射。
\end{theorem}
\begin{proof}
若 $h=\T f$，必要性显然。若 $h\notin\im\T$，则商向量空间 $F/\im\T$ 中 $h$ 的像非零。
在含该向量的一组基上定义一个取值非零的泛函，再拉回到 $F$，就得到消灭 $\im\T$ 却不消灭 $h$ 的 $\Lambda$，违反\eqref{eq:fullcompat}。

若 $\T$ 单射，对任意 $\ell\in E^*$，先在 $\im\T$ 上定义 $\Lambda(\T f)=\ell(f)$，再延拓到 $F$，故 $\T^*$ 满射。
反之，若 $0\ne f\in\ker\T$，取 $\ell(f)\ne0$，则 $\ell$ 不可能属于 $\im\T^*$。
最后，$\ker\T^*$ 正是消灭 $\im\T$ 的泛函，故其为零当且仅当 $\im\T=F$。
\end{proof}

\begin{example}[有限关系相容仍不足：标量值反例]\label{ex:noncompact}
取 $K=V=\Q,n=1$，枚举 $\Q=\{q_1,q_2,\ldots\}$。
函数 $\delta_{q_j}1$ 与常值函数 $\mathbf1$ 构成线性无关族，将其延拓为 $\Q^{\Q}$ 的 Hamel 基。
令 $\ell_j$ 为对应于基向量 $\delta_{q_j}1$ 的坐标泛函，并定义
\[
 P(f,q_j)=\ell_j(f),\qquad h(q_j)=1.
\]
任意有限子系统可由相应的有限和 $f=\sum_{j\in F}\delta_{q_j}1$ 解出，且泛函族 $(\ell_j)$ 线性无关。
但任意函数作为 Hamel 基的线性组合只用有限项，故序列 $(\ell_j(f))_{j\ge1}$ 有限支集，不可能处处等于 $1$。
所以全局方程不可解。

每个 $\ell_j$ 都非有限取值型：若它是有限个点值的线性组合，对所有 $\delta_{q_k}1$ 代入会迫使它等于 $f\mapsto f(q_j)$；然而 $\ell_j(\mathbf1)=0$ 而 $\mathbf1(q_j)=1$，矛盾。
\end{example}
这里完整对偶判据没有失效。事实上 $\im\T=\Q^{(\Q)}$，在商 $\Q^{\Q}/\Q^{(\Q)}$ 中检测常值函数的泛函提供了障碍；该泛函不可能是有限个坐标评价的线性组合。
下一节的紧性恰好使所有存在性障碍都能由有限关系检测。

\section{标量值紧线性差：关系空间控制全部解}\label{sec:scalar}
设 $W=K$，$P$ 为紧的 $K$-线性差。由上一节，存在唯一系数 $a(v,x)$，每个固定 $v$ 只有有限个 $a(v,x)$ 非零，且
\begin{equation}\label{eq:rowfinite}
 P(f,v)=\sum_{x\in I}a(v,x)f(x).
\end{equation}
令
\[
 D_I=K^{(I)}=\bigoplus_{x\in I}Ke_x,\qquad
 D_J=K^{(J)}=\bigoplus_{v\in J}Ke_v.
\]
定义行向量、关系映射与行空间
\begin{equation}\label{eq:A}
 r_v=\sum_xa(v,x)e_x,\qquad
 A:D_J\to D_I,\ A(e_v)=r_v,\qquad L=\im A.
\end{equation}
每个函数 $f:I\to K$ 唯一对应线性泛函
$\ell_f:D_I\to K$，其中 $\ell_f(e_x)=f(x)$。
因此 $K^I=D_I^*$，且 $\T=A^*$。

\begin{theorem}[存在性的四个等价条件]\label{thm:scalar}
对 $h:J\to K$，以下等价：
\begin{enumerate}
\item 存在 $f:I\to K$ 使 $P(f,v)=h(v)$ 对全部 $v$ 成立。
\item 每个有限子方程组都可解。
\item 对任意有限族 $v_1,\ldots,v_s$ 与 $c_1,\ldots,c_s\in K$，都有
\begin{equation}\label{eq:compat}
 \sum_{i=1}^s c_ir_{v_i}=0
 \quad\Longrightarrow\quad
 \sum_{i=1}^s c_ih(v_i)=0.
\end{equation}
\item $\ell_h|_{\ker A}=0$，其中 $\ell_h(e_v)=h(v)$。
\end{enumerate}
\end{theorem}
\begin{proof}
(1)$\Rightarrow$(2)显然。若有限关系 $\sum_i c_ir_{v_i}=0$ 成立，取相关有限方程的一组解，将等式乘以 $c_i$ 相加，便得到\eqref{eq:compat}，故(2)$\Rightarrow$(3)。
由于 $D_J$ 的每个元素都是有限和，(3)与(4)完全相同。

由(3)，可在 $L$ 上定义
\[
 \varphi_h\!\left(\sum_i c_ir_{v_i}\right)=\sum_i c_ih(v_i).
\]
两个表达式的差是一条有限关系，所以定义良好且线性。
取 $L$ 的一组基并延拓为 $D_I$ 的基，在新增基向量上任意指定值，便将 $\varphi_h$ 延拓为 $\ell:D_I\to K$。
令 $f(x)=\ell(e_x)$，则 $P(f,v)=\ell(r_v)=h(v)$，得到(1)。
\end{proof}

\begin{theorem}[解空间、唯一性及任意右端可解]\label{thm:scalarstructure}
记 $C=D_I/L$。若 $f_0$ 为一组解，则
\begin{equation}\label{eq:scalarstructure}
 \Sol(h)=f_0+C^*,\qquad \ker\T\cong C^*.
\end{equation}
具体地，$C^*$ 通过商映射拉回到 $D_I^*$。此外：
\begin{enumerate}
\item 可解右端的解唯一，当且仅当 $L=D_I$，等价于每个 $e_x$ 都是有限个行向量的线性组合。
\item 每个 $h$ 都可解，当且仅当 $A$ 单射，等价于行族 $(r_v)_{v\in J}$ 线性无关。
\item 每个 $h$ 恰有一解，当且仅当行族 $(r_v)$ 是 $D_I$ 的一组基。
\end{enumerate}
\end{theorem}
\begin{proof}
齐次解对应恰好消灭 $L$ 的泛函，等价于 $C$ 上的任意泛函。
非零向量空间总存在非零线性泛函：把一个非零向量纳入基并在该向量上指定值 $1$。
所以 $C^*=0$ 当且仅当 $C=0$，得到(1)。
由定理\ref{thm:scalar}，所有 $h$ 可解等价于 $D_J$ 上所有泛函都消灭 $\ker A$，亦即 $\ker A=0$，得到(2)。结合两项得到(3)。
\end{proof}

\begin{corollary}[有限矛盾证书与有限恢复公式]
不可解时，存在一条有限关系\eqref{eq:compat}，其左侧为零而右侧非零。
若解唯一，对每个 $x$ 均可选有限表达式 $e_x=\sum_i c_ir_{v_i}$，于是
\[
 f(x)=\sum_i c_i h(v_i).
\]
不同表达式给出的值因相容性而相同。
\end{corollary}

\begin{remark}[无限维数不能作形式减法]
若 $C$ 的一组基为 $B$，那么 $C^*\cong K^B$，全部自由参数是一族任意标量。
当 $|B|=d<\infty$ 时，齐次解空间维数为 $d$；当 $B$ 无限时，不可把 $K^B$ 当成直和 $K^{(B)}$，也不可写成未经解释的“未知数个数减秩”。
有限参数结构由商空间 $C$，而非由底层 $V$ 的维数直接决定。
\end{remark}

\begin{proposition}[精确序列与像的闭性]
存在精确序列
\[
 0\longrightarrow C^*\longrightarrow K^I
 \xrightarrow{\T}K^J\xrightarrow{\rho}(\ker A)^*
 \longrightarrow0,
\]
其中 $\rho(h)=\ell_h|_{\ker A}$。
赋予 $K$ 离散拓扑后，$\im\T$ 是 $K^J$ 中的闭子空间。
\end{proposition}
\begin{proof}
各处精确性来自前述延拓论证；$\rho$ 满射是因为 $\ker A$ 上的泛函可延拓到 $D_J$。
像由所有有限关系对应的方程 $\sum_i c_ih(v_i)=0$ 的交定义，每个此类方程的零集闭，故像闭。
\end{proof}

\section{有限维向量值情形的对偶判据}\label{sec:vectorfinite}
设 $0<\dim_KW=m<\infty$，并写
\[
 P(f,v)=\sum_x T_{v,x}f(x),\qquad T_{v,x}\in\End_K(W),
\]
每行仍只有有限个非零项。令
\[
 E_I=\bigoplus_{x\in I}W^*,\qquad E_J=\bigoplus_{v\in J}W^*.
\]
用 $i_x(\lambda)$ 表示位于 $x$ 坐标的 $\lambda\in W^*$，定义
\begin{equation}\label{eq:Bdual}
 B:E_J\to E_I,\qquad
 B(i_v(\lambda))=\sum_x i_x(\lambda\circ T_{v,x}).
\end{equation}
有限维性给出自然同构 $E_I^*\cong W^I$、$E_J^*\cong W^J$，且 $\T=B^*$。

\begin{theorem}[向量值存在唯一性]\label{thm:vectorfinite}
上述方程可解，当且仅当对任意有限族 $v_i\in J,\lambda_i\in W^*$，有
\begin{equation}\label{eq:veccompat}
 \sum_i\lambda_i\circ T_{v_i,x}=0\quad(\forall x\in I)
 \quad\Longrightarrow\quad
 \sum_i\lambda_i(h(v_i))=0.
\end{equation}
它也等价于每个有限子方程组可解。若可解，则
\[
 \Sol(h)=f_0+(E_I/\im B)^*.
\]
解唯一当且仅当 $B$ 满射；每个右端可解当且仅当 $B$ 单射。
\end{theorem}
\begin{proof}
在配对 $\langle f,i_x(\lambda)\rangle=\lambda(f(x))$ 下，直接计算得 $\T=B^*$。
条件\eqref{eq:veccompat}正是对应于 $h$ 的泛函消灭 $\ker B$。
同标量情形一样，先在 $\im B$ 上定义泛函，再延拓到 $E_I$；这里必须用 $W^{**}=W$ 才能把所得泛函解释为 $W$-值函数。

有限子方程可解蕴含\eqref{eq:veccompat}，因为每条关系只涉及有限个方程。
全局可解则显然蕴含有限可解。商空间、单射与满射的论证均与定理\ref{thm:scalarstructure}相同。
\end{proof}

\begin{remark}[标量系数与有限维值域是两种不同的充分条件]
若所有 $T_{v,x}=a(v,x)\id_W$，则即便 $W$ 无限维，定理\ref{thm:scalar}的证明仍可直接使用：
把 $\varphi_h:L\to W$ 延拓为 $D_I\to W$ 即可。
此时相容性是 $\sum_i c_ir_{v_i}=0\Rightarrow\sum_i c_i h(v_i)=0$，齐次解为
$\Hom_K(D_I/L,W)$。下一节的困难来自一般算子系数与无限维值域同时出现。
\end{remark}

\section{无限维值域：延拓障碍及明确反例}\label{sec:infinite}
\begin{example}[有限可解而全局不可解]\label{ex:infinite}
取 $V=\Q$、$n=1$、$K=\Q$、$W=\bigoplus_{j\ge1}\Q u_j$。
令 $\pi_j:W\to W$ 为 $\pi_j(\sum_k a_ku_k)=a_ju_1$。
在互异的参数 $v_j=j$ 处置
\[
 P(f,v_j)=\pi_j(f(0)),\qquad h(v_j)=u_1,
\]
其他参数处置 $P(f,v)=h(v)=0$。
该 $P$ 为紧的 $K$-线性差，每行甚至只取一个函数值。
任意有限子方程只要求 $f(0)$ 的有限个坐标等于 $1$，因而可解。
全部方程却要求 $f(0)=\sum_{j\ge1}u_j$，这不是 $W$ 的元素，故不可解。
\end{example}
该例说明不能在定理\ref{thm:vectorfinite}中删去有限维性。
若一般地使用 $W^*$ 作上述对偶延拓，只能先得到 $W^{**}$-值函数；
还需要保证各取值属于自然嵌入 $W\hookrightarrow W^{**}$ 的像。

\begin{theorem}[任意维值域的模同态判据]\label{thm:module}
设 $W\ne0$，$\calA=\End_K(W)$，以复合作为乘法，并把 $W$ 视为自然左 $\calA$-模。
令
\[
 \calM=\bigoplus_{x\in I}\calA e_x,\qquad
 \calN=\bigoplus_{v\in J}\calA e_v,
\]
定义左模同态
\[
 B_{\calA}:\calN\to\calM,\quad B_{\calA}(e_v)=\sum_x T_{v,x}e_x,
 \qquad\psi_h(e_v)=h(v).
\]
则 $P(f)=h$ 可解当且仅当下列两项同时成立：
\begin{enumerate}
\item $\psi_h(\ker B_{\calA})=0$，即对有限个 $a_i\in\calA$，
\[
 \sum_i a_iT_{v_i,x}=0\ (\forall x)quad\Longrightarrow\quad
 \sum_i a_i h(v_i)=0;
\]
\item 由第一项诱导的左模同态
\[
 \bar\psi_h:\im B_{\calA}\to W,\qquad
 \bar\psi_h(B_{\calA}u)=\psi_h(u)
\]
可延拓为 $\calM\to W$ 的左模同态。
\end{enumerate}
有解时，齐次解空间自然为
\[
 \Hom_{\calA}(\coker B_{\calA},W).
\]
\end{theorem}
\begin{proof}
每个 $f:I\to W$ 唯一确定左模同态
$\phi_f:\calM\to W$，其中 $\phi_f(e_x)=f(x)$。
因此
\[
 (\phi_f B_{\calA})(e_v)=\sum_xT_{v,x}f(x)=P(f,v).
\]
所求方程等价于 $\psi_h=\phi_fB_{\calA}$。
先经像分解再延拓，恰好就是两项条件。差解满足 $\phi B_{\calA}=0$，因而恰好经余核分解。
\end{proof}
这里使用\textbf{左}模，所以相容式中的乘积顺序是 $a_iT_{v_i,x}$，不能交换。
在例\ref{ex:infinite}中，任何有限算子关系都满足第一项：取相关有限子系统的一个解即可验证。
因此该例也明确展示第二项可能失败，不能无证明地删去延拓条件。

\subsection{双对偶中的形式解与真实取值障碍}
继续使用 $E_I=\bigoplus_xW^*$、$E_J=\bigoplus_vW^*$ 和式\eqref{eq:Bdual}的 $B$，
但本小节允许 $W$ 无限维。记自然嵌入为 $j:W\hookrightarrow W^{**}$。

\begin{theorem}[有限相容性的双对偶解释]\label{thm:bidual}
对任意维 $W$，以下三个条件等价：
\begin{enumerate}
\item 每个有限子方程组 $\sum_xT_{v,x}f(x)=h(v)$ 都有 $W$-值解；
\item 式\eqref{eq:veccompat}的全部有限对偶关系成立；
\item 存在 $\widehat f:I\to W^{**}$，使
\begin{equation}\label{eq:bidualsolution}
 \sum_xT_{v,x}^{**}\widehat f(x)=j(h(v))\qquad(v\in J).
\end{equation}
\end{enumerate}
原来的 $W$-值问题有解，当且仅当可在第三项中选到 $\widehat f(x)\in j(W)$ 对所有 $x$ 成立。
\end{theorem}
\begin{proof}
第一项蕴含第二项，因为每条关系只涉及有限个方程。
反过来，固定有限参数集 $F\subset J$，令 $S\subset I$ 为这些行的支集并。
有限子系统是线性映射
$T_F:W^S\to W^F$ 的一个纤维。
若 $h|_F\notin\im T_F$，线性泛函分离给出 $\Lambda\in(W^F)^*$，消灭该像却不消灭 $h|_F$。
因为 $F$ 有限，$\Lambda$ 唯一可写成 $\sum_{v\in F}\lambda_v(w_v)$，其中 $\lambda_v\in W^*$。
消灭像恰表示 $\sum_{v\in F}\lambda_v\circ T_{v,x}=0$ 对全部 $x$ 成立，这与第二项矛盾。

第二项表示由 $h$ 给出的泛函 $\ell_h(i_v\lambda)=\lambda(h(v))$ 消灭 $\ker B$。
因此它先下降为 $\im B$ 上的泛函，再延拓为 $\ell\in E_I^*$。
自然识别 $E_I^*=(W^{**})^I$，令 $\widehat f(x)(\lambda)=\ell(i_x\lambda)$，
则 $\ell B=\ell_h$ 正好是\eqref{eq:bidualsolution}。
反过来，任何这样的形式解作用到有限关系上，都给出第二项。

若各 $\widehat f(x)$ 都属于 $j(W)$，由 $j$ 的单射性和 $T^{**}j=jT$，可唯一写成 $j(f(x))$ 并得到原方程。
原方程的解经 $j$ 嵌入又显然给出这种形式解。
\end{proof}

这说明有限维性的具体作用是 $j(W)=W^{**}$。
在无限维反例\ref{ex:infinite}中，形式解允许在全部坐标上取值 $1$，但该取值不能由 $W$ 的有限支集向量实现。
模同态延拓表述与这里的双对偶表述描述的是同一个全局取值障碍。

\section{模型方程：非齐次 Cauchy、Jensen 与差分多项式}\label{sec:models}
\subsection{非齐次 Cauchy 方程的完整相容条件}
设 $V,W$ 为 $\Q$-向量空间，维数不必有限。记
\[
 Df(x,y)=f(x+y)-f(x)-f(y).
\]
\begin{theorem}[对称二余循环判据]\label{thm:cocycle}
方程 $Df=h$ 有解，当且仅当
\begin{align}
 h(x,y)&=h(y,x),\label{eq:sym}\\
 h(x,y)+h(x+y,z)&=h(y,z)+h(x,y+z)\label{eq:cocycle}
\end{align}
对全部变量成立。有解时，全部解为
\[
 f_0+\Hom_{\Q}(V,W).
\]
若 $\dim_{\Q}V=d<\infty$、$\dim_{\Q}W=m<\infty$，则解仿射空间的维数为 $dm$。
\end{theorem}
\begin{proof}
必要性由展开 $Df$ 得到。证明充分性。
在\eqref{eq:cocycle}中令 $z=0$，得 $h(x+y,0)=h(y,0)$，故 $h(x,0)$ 为常值 $c=h(0,0)$。
置 $h_0=h-c$，则 $h_0(x,0)=h_0(0,x)=0$，且仍满足两式。

在集合 $E=W\times V$ 上定义加法
\[
 (u,x)+(w,y)=(u+w+h_0(x,y),x+y).
\]
对称性给出交换律，\eqref{eq:cocycle}给出结合律；$(0,0)$ 为零元，
$(-u-h_0(x,-x),-x)$ 为 $(u,x)$ 的逆元，故 $E$ 是阿贝尔群。
对每个正整数 $k$，设
\[
 b_k(x)=\sum_{j=1}^{k-1}h_0(jx,x),
\]
则归纳得 $k(u,x)=(ku+b_k(x),kx)$。
给定 $(w,y)$，只能取 $x=y/k$，再唯一取 $u=(w-b_k(x))/k$，便有 $k(u,x)=(w,y)$。
所以 $E$ 唯一可除，因而具有唯一的 $\Q$-向量空间结构。

投影 $E\to V$ 是满射群同态，因唯一可除性也是 $\Q$-线性映射。
选取一组基的提升，得到线性截面 $s:V\to E$，写为 $s(x)=(g(x),x)$。
等式 $s(x+y)=s(x)+s(y)$ 给出
$g(x+y)=g(x)+g(y)+h_0(x,y)$，即 $Dg=h_0$。
令 $f=g-c$，因为常值 $-c$ 的缺陷是 $c$，所以 $Df=h_0+c=h$。
最后，两个解之差可加，也即 $\Q$-线性；有限维结论由基的像可任意指定得到。
\end{proof}
上述证明同样适用于任意特征零域上的向量空间，但所得自由项一般仅为底层 $\Q$-线性映射。
不能在没有额外假设时将其换成 $\Hom_K(V,W)$。
指定一组 $\Q$-基上的函数值，可以从每个非空解仿射空间中选出唯一解。

\subsection{Jensen 方程与二次方程}
\begin{proposition}[有理向量空间上的两类齐次模型]
在 $V=\Q^d,W=\Q^m$ 上：
\begin{enumerate}
\item $2f((x+y)/2)-f(x)-f(y)=0$ 的全部解为 $f(x)=A(x)+b$，其中 $A$ 为 $\Q$-线性映射。
\item $f(x+y)+f(x-y)-2f(x)-2f(y)=0$ 的全部解为 $f(x)=B(x,x)$，其中 $B:V\times V\to W$ 为对称 $\Q$-双线性映射。
\end{enumerate}
解空间维数分别为 $m(d+1)$ 与 $m d(d+1)/2$。
\end{proposition}
\begin{proof}
第一项令 $g=f-f(0)$，先取 $y=0$ 得 $2g(x/2)=g(x)$，再代回得 $g(x+y)=g(x)+g(y)$。

第二项依次取 $x=y=0$、$x=0$、$y=x$，得到 $f(0)=0$、$f(-x)=f(x)$ 和 $f(2x)=4f(x)$。
定义
\[
 B(x,y)=\frac{f(x+y)-f(x-y)}4
       =\frac{f(x+y)-f(x)-f(y)}2.
\]
后一式表明对称性。分别对 $(x+y,z)$ 与 $(x-y,z)$ 应用原式并相减，得到
$B(x+z,y)+B(x-z,y)=2B(x,y)$。
由第一项且 $B(0,y)=0$，$B$ 对第一个变量线性；对称性给出双线性。
又 $B(x,x)=f(x)$。反向代入直接成立。维数按基上的独立参数计数。
\end{proof}

\subsection{高阶混合差分与普通多项式}
定义 $\Delta_hf(x)=f(x+h)-f(x)$。
\begin{theorem}[有理有限维情形的差分分类]\label{thm:polydifference}
设 $V=\Q^d,W=\Q^s$，$m\ge0$。以下等价：
\[
 \Delta_{h_1}\cdots\Delta_{h_{m+1}}f(x)=0
 \qquad(x,h_1,\ldots,h_{m+1}\in V);
\]
$f$ 的每个坐标是总次数不超过 $m$ 的普通有理系数多项式。
解空间维数为 $s\binom{d+m}{m}$。
\end{theorem}
\begin{proof}
一个差分使多项式总次数至少下降一次，故多项式必满足方程。
反过来，只需处理标量值。在整数格上，记 $e_i$ 为标准基、$\Delta_i=\Delta_{e_i}$。
对 $\alpha\in\N^d$ 置 $c_\alpha=(\Delta_1^{\alpha_1}\cdots\Delta_d^{\alpha_d}f)(0)$。
若 $|\alpha|>m$，混合差分假设给出 $c_\alpha=0$。
由移位算子 $T_i=1+\Delta_i$ 的二项式展开，对所有 $k\in\N^d$ 有
\[
 f(k)=\sum_{|\alpha|\le m}c_\alpha
       \prod_{i=1}^d\binom{k_i}{\alpha_i}=:p(k).
\]
右边是总次数不超过 $m$ 的多项式。沿各坐标轴，方程 $\Delta_i^{m+1}f=0$ 是首尾系数非零的线性递推，已知非负整数处的值后可向负方向唯一递推，故 $f=p$ 于 $\Z^d$。

对每个正整数 $N$，把同样的论证用于格 $N^{-1}\Z^d$，得到一个总次数不超过 $m$ 的多项式 $p_N$ 与 $f$ 在该格上一致。
在共同子格 $\Z^d$ 上有 $p_N=p$。多项式若在 $\Z^d$ 上为零，逐个变量使用一元多项式只有有限个根的事实，即知其恒零。
因此 $p_N=p$。每个有理点属于某个这样的格，所以 $f=p$ 于 $\Q^d$。
总次数不超过 $m$ 的单项式数为 $\binom{d+m}{m}$，得到维数公式。
\end{proof}
\begin{remark}
当定义域换成 $\R^d$ 而不加正则性时，普通实多项式分类不成立；
例如任意不连续加性函数都满足二阶差分为零。
一般应使用底层 $\Q$-结构上的广义多项式。本文只在已证明的有理有限维范围内使用上述普通多项式结论。
\end{remark}

\section{自然性之一：值域相容与群等变}\label{sec:natural}
“自然性”不是一个无需定义的单一条件。本节和下一节分别说明值域线性相容、固定群的等变性，以及跨全部空间的拉回自然性。

\subsection{与所有值域线性自映射交换}
\begin{proposition}[值域自然性强制标量系数]\label{prop:commutant}
设 $W\ne0$，$P$ 为紧的 $K$-线性差。若对所有 $B\in\End_K(W)$ 有
\[
 P(B\circ f,v)=B(P(f,v)),
\]
则每个局部算子 $T_{v,x}$ 都是标量倍数 $a(v,x)\id_W$。
逆命题亦成立，且不要求 $W$ 有限维。
\end{proposition}
\begin{proof}
取只在 $x$ 处非零的 $f$，得到 $T_{v,x}B=BT_{v,x}$。
以下证明任何与全部自同态交换的 $T$ 都是标量。
固定 $u\ne0$，取 $\lambda\in W^*$ 使 $\lambda(u)=1$，令 $p(w)=\lambda(w)u$。
因为 $Tp=pT$，所以 $Tu=pTu\in Ku$，写为 $Tu=au$。
对任意 $w$ 取线性映射 $B$ 使 $Bu=w$，则 $Tw=TBu=BTu=aw$，故 $T=a\id_W$。
反向显然。
\end{proof}
这是一条使无限维值域重新落入标量系数理论的自然条件，见第\ref{sec:vectorfinite}节末的说明。

\subsection{固定群作用下的行向量}
设群 $G$ 作用在集合 $I,J$ 上，并在 $D_I$ 上定义 $g e_x=e_{gx}$。
\begin{proposition}[等变性与行关系]
对标量值紧线性差，条件
\begin{equation}\label{eq:equiv}
 P(f\circ g^{-1},gv)=P(f,v)\qquad(g\in G)
\end{equation}
等价于
\[
 r_{gv}=g r_v.
\]
此时 $L=\Span\{r_v\}$ 是 $K[G]$-子模，$D_I/L$ 也自然成为 $K[G]$-模。
\end{proposition}
\begin{proof}
左边展开为 $\sum_xa(gv,x)f(g^{-1}x)=\sum_ya(gv,gy)f(y)$。
对所有 $f$ 相等，当且仅当 $a(gv,gy)=a(v,y)$，这正是行向量的等变式。
行空间的稳定性随即成立。
\end{proof}
因此，等变性把任意行空间变成一个带群作用的关系空间。但固定一个空间上的群等变性，仍弱于下一节对全部线性映射的自然性。

\subsection{非齐次右端的对称性与对称解}
采用函数上的作用 $g\cdot f=f\circ g^{-1}$，输出函数同样定义作用。
\begin{proposition}[解空间的仿射群作用]\label{prop:affinesolutions}
在式\eqref{eq:equiv}成立时，
\[
 g\cdot\Sol(h)=\Sol(g\cdot h).
\]
所以右端的稳定子 $H=\{g:g\cdot h=h\}$ 作用在非空解仿射空间上。
若解唯一，则该解必为 $H$-不变函数。
\end{proposition}
\begin{proof}
等变性等价于 $\T(g\cdot f)=g\cdot\T(f)$，因此 $\T(f)=h$ 蕴含 $\T(g\cdot f)=g\cdot h$；用 $g^{-1}$ 得到集合相等。
唯一解的每个 $H$-平移仍是解，故必等于自身。
\end{proof}

\begin{proposition}[对称解的精确障碍与有限群平均]\label{prop:symmetric}
若 $f_0\in\Sol(h)$，记 $Z=\ker\T$，并对 $g\in H$ 定义
$c(g)=g\cdot f_0-f_0$。
则 $c(g)\in Z$，且
\[
 c(gk)=g\cdot c(k)+c(g).
\]
存在 $H$-不变解，当且仅当存在 $u\in Z$ 使
\[
 c(g)=u-g\cdot u\qquad(g\in H).
\]
特别地，若 $H$ 有限且 $|H|$ 在 $K$ 中非零，任意可解的 $H$-不变右端都有 $H$-不变解。
\end{proposition}
\begin{proof}
$\T(c(g))=g\cdot h-h=0$，而将 $gk\cdot f_0-f_0$ 分成
$g\cdot(k\cdot f_0-f_0)+(g\cdot f_0-f_0)$ 得到所示恒等式。
全部解为 $f_0+u$；它不变等价于 $c(g)+g\cdot u-u=0$，即所述条件。
有限群情形可以直接取
$f_{\mathrm{av}}=|H|^{-1}\sum_{g\in H}g\cdot f_0$。
线性性给出 $\T(f_{\mathrm{av}})=h$，重新排列求和项给出不变性。
\end{proof}
上述 $c$ 是解仿射空间对应的一余循环；更换基点 $f_0$ 只使它增加一个余边界。
这里的恒等式及等价条件已经给出了所需的全部内容，不依赖上同调术语。

\begin{example}[特征整除群阶时平均结论会失败]
取 $K=V=\mathbb F_p$，平移群 $H=(\mathbb F_p,+)$，并令
$P(f,v)=\sum_{x\in\mathbb F_p}f(x)$、$h(v)=1$。
这是平移等变的紧线性差，$f=\delta_0 1$ 是解。
但平移不变函数只能为常数，而常数函数的输出为 $p c=0$，所以不存在平移不变解。
\end{example}

\subsection{平移等变性与群代数}
令 $G=(V,+)$，$J=V^n$ 采用对角平移。
将 $D_V$ 识别为群代数 $K[V]$，其基记作 $[x]$，乘法为 $[x][y]=[x+y]$。
\begin{theorem}[平移等变紧算子的群代数表述]\label{thm:translation}
若 $P$ 满足平移等变性，则它由一族有限群代数元素
\[
 b_d=r_{(0,d_2,\ldots,d_n)},\qquad d\in V^{n-1},
\]
唯一确定，且
\[
 r_{(v_1,\ldots,v_n)}=[v_1]b_{(v_2-v_1,\ldots,v_n-v_1)}.
\]
行空间恰是理想
\[
 L=(b_d:d\in V^{n-1})\subseteq K[V],
\]
齐次解自然为 $(K[V]/L)^*$。
特别地，$n=1$ 时恰为固定有限模板的卷积型算子
\[
 P(f,v)=\sum_{a\in F}c_a f(v+a).
\]
\end{theorem}
\begin{proof}
对每个元组减去首坐标，应用行等变式即得公式。
所有行正是所有 $[v]b_d$，其 $K$-线性包等于所示理想。
齐次解结论来自第\ref{sec:scalar}节。$n=1$ 时只剩一个基准行 $b$，其有限支集给出卷积模板。
\end{proof}

\begin{example}[周期与高阶差分]
周期方程 $f(x+a)-f(x)=0$ 对应理想 $([a]-[0])$。
商向量空间自然具有 $V/\Z a$ 的陪集基，因此解恰为各陪集上的任意常值。
更一般地，令增广理想
\[
 \mathfrak a=\ker\bigl(K[V]\to K,\ \sum c_x[x]\mapsto\sum c_x\bigr).
\]
$\mathfrak a$ 由所有 $[h]-[0]$ 生成，
$(m+1)$ 阶混合差分的行正是
\[
 [x]\prod_{i=1}^{m+1}([h_i]-[0]),
\]
其行理想为 $\mathfrak a^{m+1}$。故广义多项式空间为 $(K[V]/\mathfrak a^{m+1})^*$。
当 $K=\Q,V=\Q^d$ 时，第\ref{sec:models}节给出它的普通多项式解释和维数。
\end{example}

\section{固定空间上的线性群与自同态自然性}\label{sec:fixednatural}
本节设 $K$ 为无限域，$0<\dim_KV=d<\infty$，$W=K$。
先仅要求 $\mathrm{GL}(V)$-等变性，不要求在不同维数的空间间相容。
给定元组 $v=(v_1,\ldots,v_n)$，定义
\[
 a_v:K^n\to V,\quad a_v(\lambda)=\sum_i\lambda_iv_i,
 \qquad N_v=\ker a_v,\quad S_v=\im a_v.
\]
商映射诱导同构 $\bar a_v:K^n/N_v\to S_v$。

\begin{lemma}[有限支集必须位于输入张成空间]\label{lem:supportspan}
若 $P$ 紧且对 $\mathrm{GL}(V)$ 等变，则
\[
 \supp r_v\subseteq S_v.
\]
\end{lemma}
\begin{proof}
元组的稳定子由逐点固定 $S_v$ 的可逆线性变换组成；等变性使 $r_v$ 在此稳定子下不变。
若 $S_v=V$，结论显然。否则取补空间 $V=S_v\oplus U$。
任意 $x=s+u\notin S_v$ 满足 $u\ne0$。
令变换在 $S_v$ 上为恒等、在 $U$ 上乘以任意 $t\in K^\times$，得到无穷多个不同点 $s+tu$。
若 $x$ 处系数非零，整个无穷轨道上的系数均相同且非零，与行的有限支集矛盾。
\end{proof}

\begin{theorem}[固定空间的线性群等变分类]\label{thm:glclassification}
上述 $\mathrm{GL}(V)$-等变紧线性差，与下列数据一一对应：
对每个余维不超过 $d$ 的子空间 $N\le K^n$，任意选取有限支集元素
\[
 b_N\in K^{(K^n/N)}.
\]
对应的行向量是
\begin{equation}\label{eq:glclassification}
 r_v=(\bar a_v)_*b_{N_v},
\end{equation}
其中 $\bar a_v$ 的值看作 $V$ 的元素。不同 $N$ 的数据可独立选取。
\end{theorem}
\begin{proof}
两元组 $v,w$ 属于同一线性群轨道，当且仅当 $N_v=N_w$。
必要性显然；充分性由 $a_v(\lambda)\mapsto a_w(\lambda)$ 给出 $S_v\to S_w$ 的同构，再延拓为 $V$ 的自同构。
由引理\ref{lem:supportspan}，每个 $r_v$ 可沿 $\bar a_v$ 唯一拉回到 $K^{(K^n/N_v)}$。
同一轨道中的拉回相同，因为行向量满足 $r_{gv}=g_*r_v$，且稳定子逐点固定其支集。
所以得到唯一的 $b_N$。
反过来，任意这些数据通过\eqref{eq:glclassification}给出行有限算子；核空间在可逆变换下不变，公式直接保证等变性。
\end{proof}

\begin{remark}[为什么要求无限域]
在有限域和有限维 $V$ 上，所有轨道均有限。
例如 $P(f,v)=\sum_{x\in V}f(x)$ 对线性群等变，但在零元组处支集可以是整个 $V$，不包含于 $S_v=0$。
因此上述支集论证及其分类不能直接用于有限域；对全部空间及全部线性映射自然的分类则没有这一限制。
\end{remark}

\begin{theorem}[对全部线性自映射相容]\label{thm:endclassification}
在定理\ref{thm:glclassification}中进一步要求，对所有 $A\in\End_K(V)$ 有
\begin{equation}\label{eq:endnaturality}
 P(f\circ A,v)=P(f,Av).
\end{equation}
则这等价于：每当 $N\subseteq N'$ 且 $\operatorname{codim}N\le d$，都有
\begin{equation}\label{eq:quotientcompat}
 b_{N'}=(q_{N,N'})_*b_N,
 \qquad q_{N,N'}:K^n/N\to K^n/N'.
\end{equation}
特别地，若 $d\ge n$，全部数据由单个有限模板 $b_0\in K^{(K^n)}$ 唯一决定，因而
\[
 P(f,v)=\sum_{\lambda\in K^n}c_\lambda f\left(\sum_i\lambda_iv_i\right).
\]
\end{theorem}
\begin{proof}
式\eqref{eq:endnaturality}对全部 $f$ 成立，当且仅当 $r_{Av}=A_*r_v$。
总有 $N_v\subseteq N_{Av}$，且在两个商空间之间 $A$ 诱导的映射正是自然商映射，故\eqref{eq:quotientcompat}足以保证该行等式。

反过来，给定 $N\subseteq N'$，分别选取核为这两个空间的元组 $v,w$。
由包含关系，规则 $a_v(\lambda)\mapsto a_w(\lambda)$ 在 $S_v$ 上良定义且线性。
将其延拓到 $V$ 的自同态 $A$，便有 $Av=w$。
此时行等式恰好给出\eqref{eq:quotientcompat}。
若 $d\ge n$，可以取 $N=0$，所以 $b_N$ 都是 $b_0$ 的商推前；唯一性和公式立即成立。
\end{proof}

当 $d<n$ 时，条件\eqref{eq:quotientcompat}仍是完整的相容数据描述，但不能在此描述中凭空使用一个未出现的满秩元组来取 $b_0$。
下一节的跨空间自然性通过允许使用通用空间 $K^n$，消除了这一维数限制。

\section{自然性之二：对全部线性映射自然的分类}\label{sec:universal}
本节研究一整族算子 $P_V:K^V\times V^n\to K$，其中 $V$ 遍历全部 $K$-向量空间。
每个 $P_V$ 均为紧的 $K$-线性差。要求对任意线性映射 $A:V\to V'$、函数 $f:V'\to K$ 和元组 $v\in V^n$，有
\begin{equation}\label{eq:functorial}
 P_V(f\circ A,v)=P_{V'}(f,Av).
\end{equation}

\begin{theorem}[线性拉回自然性的有限模板分类]\label{thm:universal}
满足\eqref{eq:functorial}的算子族与有限支集系数族 $(c_\lambda)_{\lambda\in K^n}$ 一一对应，公式为
\begin{equation}\label{eq:universal}
 P_V(f,v_1,\ldots,v_n)
 =\sum_{\lambda\in K^n}c_\lambda
 f(\lambda_1v_1+\cdots+\lambda_nv_n).
\end{equation}
此处有限支集模板不依赖 $V$ 和 $v$。
\end{theorem}
\begin{proof}
在通用空间 $U=K^n$ 中取标准元组 $(e_1,\ldots,e_n)$。
紧线性性给出唯一有限行
$r=\sum_{\lambda\in K^n}c_\lambda[\lambda]\in K^{(K^n)}$。
对任意 $v$，存在线性映射 $A:U\to V$，$A(e_i)=v_i$。
在\eqref{eq:functorial}中使用此 $A$，得到
\[
 P_V(f,v)=P_U(f\circ A,e_1,\ldots,e_n)
 =\sum_\lambda c_\lambda f(A\lambda),
\]
这就是所示公式。通用元组处的行唯一，故系数唯一。
反过来，线性映射保持所有线性组合，故任意所示模板都满足自然性。
\end{proof}

\begin{corollary}[任意维值域的自然模板]\label{cor:operatornatural}
固定任意 $K$-向量空间 $W$。若紧线性差族
$P_V:W^V\times V^n\to W$ 对全部线性映射满足式\eqref{eq:functorial}，
则存在唯一有限支集算子族 $C_\lambda\in\End_K(W)$，使
\[
 P_V(f,v)=\sum_{\lambda\in K^n}C_\lambda
 \left(f\left(\sum_i\lambda_iv_i\right)\right).
\]
若此外与 $W$ 的全部线性自映射交换，且 $W\ne0$，则每个 $C_\lambda=c_\lambda\id_W$，回到标量系数模板。
\end{corollary}
\begin{proof}
在通用空间 $K^n$ 的标准元组处应用推论\ref{cor:autolocal}，得到有限族线性局部映射。
再沿 $e_i\mapsto v_i$ 拉回，逐字重复定理\ref{thm:universal}的证明即可。
值域相容时，在通用元组处逐坐标测试，得到 $C_\lambda$ 与全部自同态交换，再用命题\ref{prop:commutant}。
\end{proof}

\begin{corollary}[额外的平移与对称性条件]\label{cor:affine}
在定理\ref{thm:universal}的假设下：
\begin{enumerate}
\item 若还要求对所有空间的平移满足\eqref{eq:equiv}，则恰好只有
\[
 \lambda_1+\cdots+\lambda_n=1
\]
处允许 $c_\lambda\ne0$。也就是说，只能对仿射组合取值。
\item 对常值函数 $f\equiv b$ 有 $P_V(f,v)=b\sum_\lambda c_\lambda$。
因此消灭所有常数等价于 $\sum_\lambda c_\lambda=0$。
\item 对元组的任意置换不变，等价于模板系数在相应坐标置换下不变。
\end{enumerate}
\end{corollary}
\begin{proof}
第一项取 $V=K^{n+1}$、$v_i=e_i$、平移量 $t=e_{n+1}$。
平移前的行再整体平移，支点为 $(\lambda,1)$；在平移后的元组上应用公式，支点为 $(\lambda,\sum_i\lambda_i)$。
比较基向量系数，非零项必须满足坐标和为 $1$。充分性由仿射组合与平移交换直接得到。
第二项直接代入。第三项在通用元组处比较有限支集系数即可。
\end{proof}

\begin{example}[Cauchy 与 Jensen 的不同自然性]
Cauchy 差的模板为
$[(1,1)]-[(1,0)]-[(0,1)]$，故对全部线性映射自然。
它\textbf{不}满足对两个输入同时平移的等变性，也不消灭非零常值函数。
Jensen 差的模板为
$2[(1/2,1/2)]-[(1,0)]-[(0,1)]$（要求 $\operatorname{char}K\ne2$），
每个支点坐标和均为 $1$、系数总和为零，因此还具有平移等变性并消灭常数。
这说明不能把仿射自然性未经检查地加到所有 Cauchy 型关系上。
\end{example}

\begin{remark}[固定空间的线性群等变性较弱]
例如在任意非零 $V$ 上定义
$P(f,v)=b(v)f(0)$，其中 $b(0)=0$、$b(v)=1$ 对 $v\ne0$。
它是紧线性差，并对 $\mathrm{GL}(V)$ 等变。
但若将 $P$ 按此公式定义于所有空间，对零线性映射 $A:V\to V'$，
在 $v\ne0$ 处左边为 $f(0)$、右边为 $0$，故不满足\eqref{eq:functorial}。
因此定理\ref{thm:universal}的“全部线性映射”假设不能换成仅有线性自同构。
\end{remark}

\section{自然模板的矩与多项式解空间}\label{sec:moments}
本节设 $\operatorname{char}K=0$，$W=K$，并使用有限模板
\[
 P_V(f,v)=\sum_{\lambda\in F}c_\lambda f\left(\sum_i\lambda_iv_i\right),
 \qquad F\subset K^n\text{ 有限}.
\]
对多重指标 $\alpha=(\alpha_1,\ldots,\alpha_n)\in\N^n$，定义模板的矩
\[
 \mu_\alpha=\sum_{\lambda\in F}c_\lambda\lambda_1^{\alpha_1}\cdots\lambda_n^{\alpha_n}.
\]
其中零次幂按 $t^0=1$ 理解，包括 $t=0$。令
\begin{equation}\label{eq:symbol}
 \sigma_m(t_1,\ldots,t_n)=\sum_{\lambda\in F}c_\lambda
 (\lambda_1t_1+\cdots+\lambda_nt_n)^m
 =\sum_{|\alpha|=m}\binom m\alpha\mu_\alpha t^\alpha.
\end{equation}

\begin{theorem}[每个齐次次数上的消灭或单射二分]\label{thm:moment}
设 $0<\dim_KV=d<\infty$，记 $H_m(V)$ 为 $V$ 上 $m$ 次齐次多项式空间。
则下列条件等价：
\begin{enumerate}
\item $P_V$ 消灭 $H_m(V)$ 中的全部多项式；
\item 对全部 $|\alpha|=m$ 有 $\mu_\alpha=0$；
\item $\sigma_m=0$。
\end{enumerate}
若这三项不成立，则 $P_V$ 在 $H_m(V)$ 上是单射。
因此对任意 $M\ge0$，
\begin{equation}\label{eq:polykernel}
 \ker P_V\cap K[V]_{\le M}
 =\bigoplus_{\substack{0\le m\le M\\\sigma_m=0}}H_m(V),
\end{equation}
其维数为 $\sum_{0\le m\le M,\ \sigma_m=0}\binom{d+m-1}{m}$。
\end{theorem}
\begin{proof}
特征零时多项式展开式\eqref{eq:symbol}的多项式系数 $\binom m\alpha$ 非零，故第二、三项等价。
当 $m\ge1$ 时，把 $p\in H_m(V)$ 写成对称 $m$-线性形式的对角值 $p(x)=B(x,\ldots,x)$；具体可取
\[
 B(x_1,\ldots,x_m)=\frac1{m!}\sum_{S\subseteq\{1,\ldots,m\}}
 (-1)^{m-|S|}p\left(\sum_{i\in S}x_i\right).
\]
对每个齐次单项式展开时，容斥只保留每个 $x_i$ 都出现一次的项，因而该式对称且多线性，并在对角线上恢复 $p$。
多线性展开得
\[
 P_V(p,v)=\sum_{|\alpha|=m}\binom m\alpha\mu_\alpha
 B(\underbrace{v_1,\ldots,v_1}_{\alpha_1},\ldots,
   \underbrace{v_n,\ldots,v_n}_{\alpha_n}).
\]
所以第二项蕴含第一项。$m=0$ 时直接用常数函数验证同一结论。

若 $\sigma_m\ne0$，由于 $K$ 无限，可取 $t\in K^n$ 使 $\sigma_m(t)\ne0$。
对每个 $x\in V$，齐次性给出
\begin{equation}\label{eq:homorecovery}
 P_V(p,t_1x,\ldots,t_nx)=\sigma_m(t)p(x).
\end{equation}
因此 $P_V(p)=0$ 强制 $p=0$，说明单射性，也说明不能消灭非零的 $H_m(V)$。

最后把 $p=\sum_{m=0}^M p_m$ 按次数分解。
若 $P_V(p)=0$，将输入元组同时乘以任意 $s\in K$，得到
$\sum_m s^mP_V(p_m,v)=0$。
无限域上一元多项式系数唯一，故每个 $P_V(p_m,v)=0$。
应用前述二分即得\eqref{eq:polykernel}，维数由齐次单项式计数给出。
\end{proof}

\begin{corollary}[非齐次多项式解的恢复与检验]\label{cor:polyrecover}
若要求 $f\in K[V]_{\le M}$，则可解右端 $h$ 必须是在元组变量中总次数不超过 $M$ 的多项式。
写其齐次分解为 $h=\sum_{m=0}^M h_m$。
对每个 $\sigma_m=0$，必须有 $h_m=0$，此时 $f_m\in H_m(V)$ 任意。
对每个 $\sigma_m\ne0$，选取 $t$ 使 $\sigma_m(t)\ne0$，唯一候选为
\[
 f_m(x)=\frac{h_m(t_1x,\ldots,t_nx)}{\sigma_m(t)}.
\]
方程可解当且仅当这些候选逐个满足 $P_V(f_m)=h_m$。
\end{corollary}
\begin{proof}
必要性与候选公式来自次数保持和\eqref{eq:homorecovery}。
若各次候选均通过代入检验，将它们与任意被消灭次数的齐次项相加便得到解。
\end{proof}
这里的最终代入检验不能省去：仅在全部变量共线的元组上恢复候选，并不自动保证对一般元组的方程成立。

\begin{proposition}[有限模板的最低非零矩有界]\label{prop:momentbound}
若模板非零，合并重合支点并删去零系数后共有 $s$ 个支点，则存在 $|\alpha|\le s-1$ 使 $\mu_\alpha\ne0$。
因此非零有限模板不能消灭所有多项式。
\end{proposition}
\begin{proof}
固定一个系数非零的支点 $\lambda_0$。
对每个其余支点 $\lambda$，选取一个坐标 $i$ 使 $\lambda_i\ne(\lambda_0)_i$，并取仿射一次多项式
$L_\lambda(z)=(z_i-\lambda_i)/((\lambda_0)_i-\lambda_i)$。
乘积 $q=\prod_{\lambda\ne\lambda_0}L_\lambda$ 次数至多 $s-1$，在 $\lambda_0$ 处为 $1$，在其余支点处为 $0$。
故 $\sum_\lambda c_\lambda q(\lambda)=c_{\lambda_0}\ne0$。
把 $q$ 展开为单项式线性组合，便有至少一个相应的矩非零。
\end{proof}

\begin{example}[Cauchy、Jensen 与差分的次数结构]
Cauchy 差的符号为 $\sigma_m(t,u)=(t+u)^m-t^m-u^m$。
其中 $\sigma_0=-1$、$\sigma_1=0$，而 $m\ge2$ 时混合项非零，所以多项式核恰为一次齐次多项式。
Jensen 差的符号为 $2((t+u)/2)^m-t^m-u^m$；它恰在 $m=0,1$ 时为零，所以多项式核恰为仿射函数。

对 $r\ge1$，重复步长差分
\[
 \Delta_h^r f(x)=\sum_{k=0}^r(-1)^{r-k}\binom rk f(x+kh)
\]
有 $r+1$ 个模板支点 $(1,k)$。
其次数小于 $r$ 的矩全为零，而
$\mu_{(m-r,r)}=\sum_k(-1)^{r-k}\binom rk k^r=r!\ne0$ 对每个 $m\ge r$ 成立。
该恒等式来自一元 $r$ 阶差分将 $t^r$ 化为常数 $r!$。
故其多项式核恰为次数小于 $r$ 的多项式，且命题\ref{prop:momentbound}的界可达到。
\end{example}
上述结论只是在多项式函数类内部的分类。对任意函数求全局核时，不得据此排除不连续加性函数等非多项式解。

\section{一元自然算子：乘法群代数与完整解判据}\label{sec:dilation}
现在取 $n=1$、$V\ne0$、$W=K$，考虑自然模板
\begin{equation}\label{eq:dilation}
 P(f,v)=\sum_{\lambda\in K}c_\lambda f(\lambda v),
\end{equation}
系数有限支集。本节的群代数判据适用于任意域。
记
\[
 s=\sum_{\lambda\in K}c_\lambda,\qquad
 \mathcal R=K[K^\times],\qquad
 b=\sum_{\lambda\in K^\times}c_\lambda[\lambda]\in\mathcal R.
\]
每条一维子空间 $L\subseteq V$ 选非零代表 $u_L$。

\begin{theorem}[按射影直线分解非齐次问题]\label{thm:dilationsolve}
方程 $P(f)=h$ 有解，当且仅当可选择 $z\in K$，使 $sz=h(0)$，并且对每条 $L$，函数
\[
 k_L(t)=h(tu_L)-c_0z\qquad(t\in K^\times)
\]
对应的线性泛函 $\ell_{k_L}:\mathcal R\to K$ 消灭
\[
 \operatorname{Ann}_{\mathcal R}(b)=\{a\in\mathcal R:ba=0\}.
\]
固定这样的 $z=f(0)$ 后，各条 $L$ 上的解相互独立；每一条上的自由解空间为 $(\mathcal R/b\mathcal R)^*$。
\end{theorem}
\begin{proof}
原点方程为 $sf(0)=h(0)$。
对 $v=tu_L\ne0$，置 $g_L(t)=f(tu_L)$，方程变为
\[
 \sum_{\lambda\ne0}c_\lambda g_L(\lambda t)=k_L(t).
\]
将 $g_L$ 识别为 $\mathcal R$ 上的线性泛函，此算子就是乘法映射
$M_b:\mathcal R\to\mathcal R,\ a\mapsto ba$ 的代数对偶。
第\ref{sec:scalar}节的关系空间判据给出相容条件 $\ell_{k_L}|_{\ker M_b}=0$，自由解为余核的对偶。
所有非零向量恰好属于一条这样的 $L$，且式\eqref{eq:dilation}不混合不同 $L$，所以这些解可以独立拼接。
\end{proof}

\begin{corollary}[一元自然算子的整体存在唯一性]\label{cor:dilationbij}
对式\eqref{eq:dilation}有：
\begin{enumerate}
\item 算子 $\T:K^V\to K^V$ 单射，当且仅当 $s\ne0$ 且 $b$ 为 $\mathcal R$ 中的单位；在这种情形下它自动为双射。
\item 每个 $h$ 都可解，当且仅当 $s\ne0$ 且 $M_b$ 单射，即 $b$ 的零化子为零。
\item 因此，至多一解与对任意右端恰有一解，对这类一元自然算子等价。
\end{enumerate}
\end{corollary}
\begin{proof}
若 $s=0$，所有常值函数都在核中，且非零 $h(0)$ 无法实现，所以单射和满射都不成立。
设 $s\ne0$。齐次解的原点值必须为零；其余自由项为各条 $L$ 上的 $(\mathcal R/b\mathcal R)^*$。
因此核为零当且仅当 $b\mathcal R=\mathcal R$，这又等价于 $b$ 为单位；此时每条直线上的乘法映射可逆，所以所有右端均可唯一求解。
任意右端可解则等价于 $M_b^*$ 满射，按代数对偶判据等价于 $M_b$ 单射。
\end{proof}

\begin{example}[反射模板的显式分类]
设 $\operatorname{char}K\ne2$，考虑 $P(f,v)=f(v)-a f(-v)$。
若 $a^2\ne1$，则任意右端均有唯一解
\[
 f(v)=\frac{h(v)+a h(-v)}{1-a^2}.
\]
这由原方程与在 $-v$ 处的方程联立得到，代入也可直接验证。
若 $a=1$，可解当且仅当 $h(-v)=-h(v)$，全部解为 $h/2$ 加任意偶函数。
若 $a=-1$，可解当且仅当 $h(-v)=h(v)$，全部解为 $h/2$ 加任意奇函数。
\end{example}

\begin{example}[任意右端可解而不唯一]
在特征零域上取 $P(f,v)=f(2v)-2f(v)$。
原点处强制 $f(0)=-h(0)$。
在非零向量集合上按 $v\sim2^k v$ 分成轨道，每条轨道选一个代表 $u$，任意指定 $f(u)$，然后用
\[
 f(2^{k+1}u)=2f(2^ku)+h(2^ku)
\]
向正、负两个方向递推。
特征零保证轨道没有周期，因此得到全部解。任意 $h$ 均可解，每条非零轨道各有一个自由参数。
\end{example}

\section{加法 Cauchy 相关集：二次型与复不可约性}\label{sec:geometry}
从本节起，$D_f(x,y)=f(x+y)-f(x)-f(y)$，$X_D(f)=Z(D_f)$。
加上线性函数不改变 $D_f$；乘以非零标量不改变 $X_D(f)$。

\subsection{复二次多项式的秩分类}
\begin{theorem}[二次型相关集]\label{thm:quadratic}
设 $f:\C^d\to\C$ 为
$f(z)=z^{\mathsf T}Az+b^{\mathsf T}z+c$，其中 $A=A^{\mathsf T}$，$r=\rank A$。
则
\[
 D_f(x,y)=2x^{\mathsf T}Ay-c.
\]
作可逆线性变换后，相关集同构于
\[
 \left\{(u,v)\in\C^r\times\C^r:\sum_{i=1}^r u_iv_i=c/2\right\}
 \times\C^{2(d-r)}.
\]
若 $r=0$，它在 $c=0$ 时为全部空间，否则为空；若 $r\ge1$，它是非空超曲面，并有：
\begin{enumerate}
\item $c\ne0$ 时光滑且不可约；
\item $c=0,r=1$ 时为两个不同超平面之并；
\item $c=0,r\ge2$ 时不可约，奇点集为 $\ker A\times\ker A$。
\end{enumerate}
\end{theorem}
\begin{proof}
展开二次式得第一式；双线性配对的秩标准形给出第二式。
若 $c\ne0$，关于 $u$ 的系数是 $v_1,\ldots,v_r$，常数项是非零常数 $-c/2$，它们在 $\C[v]$ 中的最大公因子为 $1$。
该多项式在 $\C(v)[u]$ 中是非零一次式，故不可约，由 Gauss 引理知在原多项式环中亦不可约。
若 $c=0,r\ge2$，系数 $v_1,\ldots,v_r$ 的最大公因子仍为 $1$，同理不可约。
若 $r=1,c=0$，方程就是 $u_1v_1=0$。
梯度为 $(2Ay,2Ax)$，因此只有 $Ax=Ay=0$ 才可能为奇点；这些点在 $c=0$ 时全部落在零集上，在 $c\ne0$ 时均不在零集上。
\end{proof}

\subsection{复一元多项式的不可约性完整判据}
\begin{theorem}[非空不可约相关集的分类]\label{thm:irreducible}
设 $f\in\C[z]$。$X_D(f)$ 是非空不可约代数集，当且仅当以下之一成立：
\begin{enumerate}
\item $f(z)=az$，此时 $X_D(f)=\C^2$；
\item $\deg f\ge2$ 且 $f(0)\ne0$，此时 $D_f$ 本身不可约，$X_D(f)$ 是不可约曲线。
\end{enumerate}
若 $\deg f\le1$ 且 $f(0)\ne0$，相关集为空；若 $\deg f\ge2$ 且 $f(0)=0$，相关集可约并包含两坐标轴。
\end{theorem}
\begin{proof}
一次及常数情形直接计算 $D_f=-f(0)$。
设 $d=\deg f\ge2$。若 $f(0)=0$，则 $D_f(x,0)=D_f(0,y)=0$，所以 $xy\mid D_f$。
两坐标轴都是其不同不可约分支，因此相关集可约。

若 $f(0)=c\ne0$，假设 $D_f=GH$ 为两个非常数多项式之积。
限制到 $y=0$，得到 $G(x,0)H(x,0)=-c$，因此两者均为非零常数；限制到 $x=0$ 亦然。
同一因子的两条轴上常数相同，均等于该因子在原点的值，故
\[
 G=G(0,0)+xyG_1,\qquad H=H(0,0)+xyH_1.
\]
因为两个因子非常数，其最高齐次部分均被 $xy$ 整除，故 $D_f$ 的最高齐次部分应被 $x^2y^2$ 整除。
但该最高齐次部分是
\[
 a_d\bigl((x+y)^d-x^d-y^d\bigr),
\]
其中 $x^{d-1}y$ 的系数为 $d a_d\ne0$，矛盾。故 $D_f$ 不可约。
非恒定复二元多项式有零点，所以相关集非空。
\end{proof}

\begin{example}[次数七的重因子]
对 $f(z)=z^7$，
\[
 D_f(x,y)=7xy(x+y)(x^2+xy+y^2)^2.
\]
因此在讨论作为集合的相关集时，必须区分定义多项式与其无重因子部分。
\end{example}

\section{哪些复代数集能成为加法 Cauchy 相关集}\label{sec:realization}
\begin{theorem}[缺陷多项式的微分判据]\label{thm:diffcriterion}
给定 $F\in\C[x,y]$，存在 $f\in\C[z]$ 使 $F=D_f$，当且仅当
\begin{align}
 F(x,0)=F(0,y)&=F(0,0),\label{eq:axes}\\
 F_{xxy}&=F_{xyy}.\label{eq:pde}
\end{align}
$f$ 在加上线性项 $az$ 的意义下唯一，且可由 $F$ 显式恢复。
\end{theorem}
\begin{proof}
若 $F=D_f$，则轴上均为 $-f(0)$，且 $F_{xy}=f''(x+y)$，故必要。
反过来，令 $H=F_{xy}$。\eqref{eq:pde}给出 $H_x=H_y$。
在变量 $s=x+y,t=x-y$ 下，这等价于 $H$ 不依赖 $t$，即 $H=q(x+y)$，其中 $q$ 为一元多项式。
取唯一满足 $f_0''=q$、$f_0(0)=f_0'(0)=0$ 的多项式 $f_0$。
则 $(F-D_{f_0})_{xy}=0$，故 $F-D_{f_0}=a(x)+b(y)$。
它在两轴上的限制均为常数 $k=F(0,0)$，从而自身恒等于 $k$。
令 $f=f_0-k$ 即有 $D_f=D_{f_0}+k=F$。
若 $D_f=D_g$，则 $(f-g)''=0$，再由缺陷为零知常数项也为零，因此 $f-g=az$。
\end{proof}

\begin{lemma}[不可约因子的重数最多为二]\label{lem:multiplicity}
若 $\deg f\ge2$，则 $D_f$ 的每个不可约因子重数至多为 $2$。
\end{lemma}
\begin{proof}
若不可约 $q$ 满足 $q^3\mid D_f$，二次微分后仍有 $q\mid(D_f)_{xy}=f''(x+y)$。
在 $\C$ 上，$f''(x+y)$ 的不可约因子只能是 $x+y-a$，故 $q$ 与此线性式相伴。
然而若 $(x+y-a)^2\mid D_f$，则 $(D_f)_x=f'(x+y)-f'(x)$ 限制在 $y=a-x$ 上应恒零，给出 $f'(a)-f'(x)\equiv0$。
这迫使 $f'$ 为常数，与 $\deg f\ge2$ 矛盾。
\end{proof}

\begin{theorem}[代数集的可实现性判据]\label{thm:realization}
空集和 $\C^2$ 都可实现。对非空真代数集 $Y\subsetneq\C^2$，存在多项式 $f$ 使 $Y=X_D(f)$，当且仅当：
\begin{enumerate}
\item $Y$ 为纯一维代数集，写成无重因子方程 $Y=Z(q)$，其中 $q=q_1\cdots q_s$、$q_j$ 两两不相伴且不可约；
\item 可选择 $\varepsilon_j\in\{1,2\}$，使
\[
 F=\prod_{j=1}^s q_j^{\varepsilon_j}
\]
满足\eqref{eq:axes}和\eqref{eq:pde}。
\end{enumerate}
对非空真\textbf{不可约}曲线 $Y=Z(q)$，判据进一步简化为：
$Y$ 与两坐标轴均不相交，并且 $q_{xxy}=q_{xyy}$。
\end{theorem}
\begin{proof}
取 $f\equiv1$ 得空集，取 $f=0$ 得全部空间。
非空真相关集由非恒定 $D_f$ 定义，故为纯一维。
复多项式环中唯一分解和零点定理表明，若 $Z(D_f)=Z(q)$，则
$D_f=c\prod_jq_j^{\varepsilon_j}$，$c\ne0$、$\varepsilon_j\ge1$。
由引理\ref{lem:multiplicity}，每个指数不超过二；除以常数 $c$ 不改变两项线性判据，故得到必要性。
反之，定理\ref{thm:diffcriterion}给出 $F=D_f$，其零集恰为 $Y$。

若 $Y$ 不可约，定理\ref{thm:irreducible}说明实现它的 $D_f$ 本身不可约且 $f(0)\ne0$，从而 $q$ 与 $D_f$ 相伴，且曲线避开两轴。
反过来，避开 $y=0$ 表明一元多项式 $q(x,0)$ 无根，故为非零常数；另一条轴同理，两个常数均为 $q(0,0)$。
于是轴条件成立，再用微分判据即得实现。
\end{proof}
判据使用无重因子方程的任意非零常数倍均等价。
例 $xy-1=0$ 可以实现，例 $x+y-1=0$ 因与轴相交不能作为非空不可约的 $X_D(f)$。

\section{实二次、三次相关集及两种光滑要求}\label{sec:real}
\subsection{二次与三次多项式}
对 $f(t)=at^2+bt+c$，有 $D_f=2axy-c$。
当 $a\ne0,c\ne0$ 时为光滑实双曲线；$a\ne0,c=0$ 时为坐标轴之并；
$a=0$ 时按 $c=0$ 或 $c\ne0$ 分别为全部平面或空集。

\begin{proposition}[三次曲线的奇点和可约性]\label{prop:cubic}
设 $f(t)=at^3+bt^2+ct+d$，$a\ne0$，令 $A=3a$、$B=2b$。
则
\[
 D_f(x,y)=xy(A(x+y)+B)-d.
\]
当 $d=0$ 时，它是三条互异直线 $x=0,y=0,A(x+y)+B=0$ 的并。
当 $d\ne0$ 时，它在 $\C$ 上不可约，且有奇点当且仅当
\[
 B\ne0,\qquad d=\frac{B^3}{27A^2}
 =\frac{8b^3}{243a^2}.
\]
该情形唯一奇点为 $x=y=-B/(3A)=-2b/(9a)$，它是实孤立双点。
其他 $d\ne0$ 情形均为光滑的几何不可约实代数曲线。
\end{proposition}
\begin{proof}
直接展开得公式；$d=0$ 的分解显然，$d\ne0$ 的不可约性来自定理\ref{thm:irreducible}。
对 $F=xy(A(x+y)+B)-d$ 有
\[
 F_x=y(2Ax+Ay+B),\qquad F_y=x(Ax+2Ay+B).
\]
$F=0,d\ne0$ 时 $x,y$ 均非零，故梯度为零迫使
$2Ax+Ay+B=Ax+2Ay+B=0$，解得 $x=y=-B/(3A)$。
代回得到 $d=B^3/(27A^2)$，并因 $d\ne0$ 得 $B\ne0$。
此时 Hessian 矩阵为
\[
 -\frac{B}{3}\begin{pmatrix}2&1\\1&2\end{pmatrix},
\]
它非退化且实定号，故局部零集只有该孤立点，而复切锥为两条不同的共轭直线。
\end{proof}

\subsection{一般实多项式：实点正则与代数簇光滑}
设 $f\in\R[t]$ 且 $\deg f\ge2,f(0)\ne0$。
由定理\ref{thm:irreducible}，$F=D_f$ 在 $\C[x,y]$ 中不可约。
因此它在 $\R$ 上几何不可约，且已是约化方程。
\begin{proposition}[可直接检验的光滑判据]\label{prop:smooth}
曲线的奇点恰由下列方程组描述：
\begin{equation}\label{eq:singularsystem}
 \begin{cases}
 f(x+y)=f(x)+f(y),\\
 f'(x+y)=f'(x)=f'(y).
 \end{cases}
\end{equation}
所有\textbf{实点}均代数正则，当且仅当该方程组无实解。
曲线 $\operatorname{Spec}\R[x,y]/(F)$ 在 $\R$ 上光滑，当且仅当该方程组无复解。
前一条件使非空实点集成为光滑的一维实流形，后一条件更强。
\end{proposition}
\begin{proof}
$F_x=f'(x+y)-f'(x)$、$F_y=f'(x+y)-f'(y)$，故方程组恰为 $F=F_x=F_y=0$。
约化平面超曲面在点处正则等价于梯度非零；实点处由隐函数定理即为光滑实曲线。
特征零域是完美域，有限型超曲面在 $\R$ 上光滑等价于复化后无几何奇点，故应检查所有复解。
\end{proof}
这里“实点代数正则”是精确的梯度条件。如果“光滑”仅指点集作为嵌入流形的光滑性，一般不能忽略约化和实代数结构后直接断言与任意所给方程的梯度条件等价。

\begin{example}[实点全部光滑但代数簇不光滑]\label{ex:smoothcounter}
取
\[
 f(t)=(t^2+1)^3+1=t^6+3t^4+3t^2+2.
\]
$f(0)=2$，故 $F=D_f$ 在 $\C$ 上不可约。
对实数 $t$，
\[
 f''(t)=30t^4+36t^2+6>0,
\]
所以 $f'$ 严格递增。若实数 $x,y$ 满足梯度为零，则
$f'(x+y)=f'(x)=f'(y)$ 强制 $x+y=x=y$，即 $x=y=0$。
但 $F(0,0)=-2\ne0$，所以曲线无实奇点。
它有实点，因为 $F(0,0)<0$ 而 $F(t,t)>0$ 对充分大的实数 $t$ 成立。

另一方面，$f(i)=f(-i)=1$、$f(0)=2$，且 $f'(i)=f'(-i)=f'(0)=0$。
所以 $(i,-i)$ 与 $(-i,i)$ 都是复奇点。于是实点集是一条非空光滑实曲线，但对应的实代数簇不光滑。
\end{example}

\section{乘法 Cauchy 关系：函数解与仿射族刚性}\label{sec:multiplicative}
本节使用不同的缺陷
\[
 M_f(x,y)=f(x+y)-f(x)f(y),\qquad X_M(f)=Z(M_f).
\]
它一般不是关于 $f$ 的线性差，也不是第\ref{sec:compact}节的紧算子：乘积项带有两个函数值之间的交互。

\subsection{全域解与保持相关集的变换}
\begin{proposition}[全域解的代数形式]
设 $G$ 为阿贝尔群，$S$ 为含幺交换环。方程 $M_f=0$ 的全部解是零函数，以及如下函数：
选非零幂等元 $e\in S$，再选群同态
$\chi:G\to(eS)^\times$，并将其视为 $S$-值函数。
\end{proposition}
\begin{proof}
令 $e=f(0)$，则 $e^2=e$、$f(x)=ef(x)$、$f(x)f(-x)=e$。
故 $e=0$ 时 $f=0$；否则 $f(x)$ 都是以 $e$ 为单位元的环 $eS$ 中的单位，原方程就是群同态条件。反向显然。
\end{proof}
当 $G$ 为实向量空间、$S=\R$ 时，非零解满足 $f(x)=f(x/2)^2>0$，故全部解为 $0$ 或 $e^A$，$A$ 为任意加性函数。
当 $G=S=\Q$ 时，非零解只有 $1$：对任意素数 $p$ 与正整数 $m$，
$v_p(f(x))=m v_p(f(x/m))$，所以每个赋值为零，结合正性得 $f(x)=1$。

若 $\chi:G\to S^\times$ 为群同态，则
\[
 M_{\chi f}(x,y)=\chi(x+y)M_f(x,y),
\]
从而相关集不变。域值且 $f$ 处处非零时，还可通分得到
\[
 M_{1/f}(x,y)=-\frac{M_f(x,y)}{f(x+y)f(x)f(y)},
\]
所以 $X_M(1/f)=X_M(f)$。

\subsection{实参数同域族的子群结构与符号律}
令 $f,g:G\to\R$，$d=f-g$，$h_\lambda=g+\lambda d$。
\begin{theorem}[实仿射族刚性]\label{thm:realrigidity}
若 $X_M(h_\lambda)$ 与全部实参数 $\lambda$ 无关，且 $f\ne g$，则
\[
 f(0)=g(0)=1,
\]
$E=\{x:d(x)=0\}$ 为真加法子群，
$\chi=f|_E=g|_E:E\to\R^\times$ 为群同态，且
\[
 d(x+e)=\chi(e)d(x)\quad(e\in E),\qquad
 d(x)d(-x)<0\quad(x\notin E).
\]
因此在同域假设下，$f=g$ 当且仅当 $f-g$ 不变号。
\end{theorem}
\begin{proof}
固定 $x,y$，有参数多项式
\begin{align*}
 M_{h_\lambda}(x,y)={}&M_g(x,y)
 +\lambda\bigl(d(x+y)-g(x)d(y)-g(y)d(x)\bigr)\\
 &-\lambda^2d(x)d(y).
\end{align*}
同域条件等价于每个这样的多项式恒零或没有实根。
在 $(0,0)$ 处，若 $d(0)\ne0$，则 $h_\lambda(0)-h_\lambda(0)^2$ 是有根的非零二次式，矛盾。
故 $d(0)=0$。记 $f(0)=g(0)=c$，取 $d(x)\ne0$，则在 $(x,0)$ 处的多项式为
$(1-c)(g(x)+\lambda d(x))$，所以 $c=1$。

固定 $e\in E$，参数多项式的一次项必须为零，因非恒定实一次式一定有根。
于是 $d(x+e)=g(e)d(x)$。
若 $g(e)=0$，平移的满射性迫使 $d=0$，矛盾，所以 $g(e)\ne0$。
这既给出 $E$ 对加法封闭，也由 $x=-e$ 给出对逆元封闭。
取一个 $d(u)\ne0$，比较两次平移与一次平移，即得
$g(e_1+e_2)d(u)=g(e_1)g(e_2)d(u)$，约去 $d(u)$ 得到 $\chi$ 的乘法性。

若 $x\notin E$，则 $-x\notin E$。
在 $(x,-x)$ 处的多项式为
$1-(g(x)+\lambda d(x))(g(-x)+\lambda d(-x))$。
它在 $\lambda=-g(x)/d(x)$ 处等于 $1$，故非恒零又无实根，因而处处为正。
其二次项系数 $-d(x)d(-x)$ 非零，故必须为正。这就给出符号律及最后的等价结论。
\end{proof}

\begin{theorem}[复参数完全刚性]
设 $f,g:G\to\C$。则
\[
 X_M(\lambda f+(1-\lambda)g)\text{ 与 }\lambda\in\C\text{ 无关}
 \quad\Longleftrightarrow\quad f=g.
\]
\end{theorem}
\begin{proof}
若存在 $x$ 使 $d(x)\ne0$，在 $(x,x)$ 处参数多项式的二次项为 $-d(x)^2\ne0$。
它有复根，所以同域条件强迫它在每个参数处为零，与非零二次项矛盾。
逆向显然。
\end{proof}
实参数不能换成复参数后仍保留非平凡例子。
例如 $h_\lambda(x)=\cos x+\lambda\sin x$ 满足
$M_{h_\lambda}(x,y)=-(1+\lambda^2)\sin x\sin y$，
所有实参数同域，而 $\lambda=\pm i$ 时相关集变为全域。

\section{乘法相关域的实复拓扑差异}\label{sec:topology}
这里“区域”专指欧氏拓扑中的非空开连通集；不假设原函数连续。

\begin{lemma}[区间上的局部乘法分离]\label{lem:pexider}
设 $U,T\subseteq\R$ 为非空开区间，且 $H(s+t)=F(s)G(t)$ 对 $s\in U,t\in T$ 成立。
则或者 $F\equiv0$ 或 $G\equiv0$ 且 $H\equiv0$；或者存在唯一加性函数 $A:\R\to\R$ 及非零实常数 $a,b$，使
\[
 F(s)=ae^{A(s)},\quad G(t)=be^{A(t)},\quad H(u)=ab e^{A(u)}.
\]
\end{lemma}
\begin{proof}
假设 $F,G$ 均不恒零，取 $F(s_0)G(t_0)\ne0$。
小 $h$ 满足 $F(s_0+h)G(t_0-h)=H(s_0+t_0)\ne0$，故两函数在这些点附近非零。
定义 $c(h)=G(t_0+h)/G(t_0)$。比较同一和 $s+h+t_0=s+(t_0+h)$，得
\[
 F(s+h)=c(h)F(s)
\]
对小 $h$ 及 $s,s+h\in U$ 成立。
在 $s_0$ 附近重复平移，得到局部乘法律 $c(h+k)=c(h)c(k)$，
故 $c(h)=c(h/2)^2>0$，$a_0(h)=\log c(h)$ 局部可加。

对任意 $t\in\R$，取充分大正整数 $m$，令 $A(t)=m a_0(t/m)$。
不同的 $m$ 可用公倍细分比较，局部加性保证定义与选择无关；同时细分两个数及其和，得到 $A$ 全局可加。
把区间内的任意位移分成有限个足够小的等长段，即得
$F(s)=F(s_0)e^{A(s-s_0)}$。
交换两个变量并比较小位移，得到 $G$ 具有同一 $A$，随后得到 $H$ 的形式。
若两种表示给出同一 $F$，其加性指数之差在 $U-U$ 上为零；该集含零邻域，故加性迫使该差处处为零。常数亦唯一。
\end{proof}

\begin{theorem}[实区域必为星形集]\label{thm:topreal}
若 $f:\R\to\R$ 且 $X=X_M(f)$ 为区域，则 $X$ 为星形集，故可缩且 $\pi_1(X)=0$。
\end{theorem}
\begin{proof}
在 $X$ 内每个开矩形上应用引理\ref{lem:pexider}，知乘积 $f(x)f(y)$ 或处处为零，或处处非零。
于是 $X_0=\{(x,y)\in X:f(x)f(y)=0\}$ 及其补集都在 $X$ 中开。
连通性迫使其中一个为空。

若乘积处处非零，则 $f$ 无零点，因为 $f(z)=0$ 会使 $(0,z)\in X_0$。
记 $I=\operatorname{pr}_1X=\operatorname{pr}_2X$、$J=\{x+y:(x,y)\in X\}$，二者均为开区间。
用相交开矩形链连接 $X$ 的任意两点；相邻矩形的坐标区间相交，由局部表示唯一性，同一加性函数与常数可以拼接。
因此 $f=ae^A$ 于第一投影、$f=be^A$ 于第二投影。
两个投影同为 $I$，故 $a=b$；并有 $f=a^2e^A$ 于 $J$。
从而
\[
 X=(I\times I)\cap\{(x,y):x+y\in J\},
\]
这是一凸集。

若乘积处处为零，令 $Z=f^{-1}(0)$。
一方面 $x+y\in Z$ 对 $(x,y)\in X$ 成立；另一方面 $z\in Z$ 时 $(0,z)\in X$，所以 $Z$ 恰为和映射的像，因而是开区间。
此时
\[
 X=\{x\in Z,x+y\in Z\}\ \cup\ \{y\in Z,x+y\in Z\}.
\]
这两个集合均非空、开且凸，连通性迫使二者相交。
任取交中的一点为星心，即见并集星形。
星形集的直线收缩给出可缩性及基本群结论。
\end{proof}

\begin{theorem}[复区域可实现任意有限秩自由群]
对每个 $r\ge0$，存在只取 $2,3,4$ 三个值的函数 $f_r:\C\to\C$，使 $X_M(f_r)$ 为区域且
$\pi_1(X_M(f_r))\cong F_r$，其中 $F_r$ 为秩 $r$ 的自由群。
\end{theorem}
\begin{proof}
取圆盘 $B=D(3,1)$。
在 $D(6,1)$ 内删去 $r$ 个两两不交且完全位于内部的小闭圆盘，所得区域记作 $A$；
可使所删圆盘均在 $D(6,3/4)$ 内。于是 $A\cap B=\varnothing$ 且 $A/2\subset B$。
定义 $f_r=2$ 于 $B$、$f_r=4$ 于 $A$、其余处为 $3$。
值集 $\{2,3,4\}$ 中唯一仍属于该集合的二元乘积是 $2\cdot2=4$，故
\[
 X_M(f_r)=\{(x,y)\in B\times B:x+y\in A\}.
\]
它是开集。令 $z=x+y$，同伦
\[
 H_t(x,y)=\bigl((1-t)x+tz/2,(1-t)y+tz/2\bigr)
\]
保持和为 $z$，由 $B$ 凸且 $z/2\in B$，始终留在该集合内。
它给出到 $\{(z/2,z/2):z\in A\}\cong A$ 的强变形收缩。
因此该集合路径连通；而有 $r$ 个洞的圆盘与 $r$ 个圆的楔和同伦等价，基本群为 $F_r$。
\end{proof}

\section{本稿范围、核查要点与后续问题}\label{sec:scope}
本稿以2026年10月8日“讨论线性差解的存在性”和“紧算子等价刻画”的问题为基础，
纳入同日“讨论Cauchy关系的P面积”及已有乘法 Cauchy 草稿中直接相关的结果。
基础对象与初等方程沿用既有 \texttt{P\_operator\_draft.tex}，而“线性差”与“加性差”统一按10月9日修订后的定义使用；
乘法关系部分与 \texttt{multiplicative\_Cauchy\_draft.tex} 的定义一致。
第\ref{sec:natural}--\ref{sec:dilation}节将自然性方向整理为明确假设下可独立验证的命题。

\medskip
\noindent\textbf{本次增补的依据与范围。}
已确认10月9日讨论修订了线性差的定义，并继续研究紧线性差对 $V$ 的自然性。
该次后续讨论的完整答复原文未能取得，因此本次新增的完整代数对偶判据、双对偶解释、固定空间自然性分类、模板矩及一元伸缩方程均作为沿此方向独立补齐证明的内容；本稿不声称已逐条归档未取得的会话内容。

以下边界在后续使用时应保留：
\begin{enumerate}
\item 紧性不仅是有限依赖，还排除不同函数值之间的交互项；可加情形下后一条件才自动成立。
\item 加性差只自动给出素域线性性；线性差按定义是 $K$-线性的。算子对函数变量为 $K$-线性，并不意味着解函数对定义域变量为 $K$-线性；非齐次 Cauchy 自由项一般仅为 $\Q$-线性的。
\item 一般算子系数下，有限相容性推出全局可解需要有限维值域；无限维时保留模同态延拓障碍。标量系数是另一种可消除障碍的条件。
\item 无限维解空间由代数对偶或 Hom 空间描述，不以未经定义的无限基数相减描述。
\item 平移等变、固定空间的线性群等变、全部线性映射的自然性互不等同；每次使用分类都要检查具体假设。
\item 对称右端不自动保证存在对称解；固定空间线性群的支集分类需要无限域，而模板矩的次数二分在本文要求特征零。
\item 实点光滑与整个实代数簇光滑有区别；后者须排除复化后的奇点。
\end{enumerate}

仍可继续研究的方向包括：有限域上固定空间的等变分类，以及 $\dim V<n$ 时式\eqref{eq:quotientcompat}中相容数据的进一步简化；
对自然有限模板施加对称性、归一化后，任意函数解与多项式解之间的差异及关系模的更细结构；
给定 $h$ 与指定真子集 $Y\subset J$，研究 $X_{P-h}(f)=Y$ 的实现问题；
以及在连续、可测或多项式函数类内部重新讨论延拓与相容性。
这些额外约束下的完整分类不由本文的任意函数理论自动给出，本稿不将其列作已完成结论。

\end{document}
